\scanpage{25}\mspage{24}
Notation $F\preccurlyeq G$. C'est encore une relation d'ordre sur $\unsure{\mathcal F|\mathcal J}$.

\uline{NB} Se donner dans $L$ un couple $F\subset G$, avec $F$ raffinant $G$, revient au $=$ que de se donner
\diagram{\begin{tikzcd}[ampersand replacement=\&,column sep=large,row sep=large]
S\;(=|F|)\arrow[r,hook,"i"]\arrow[d,"p","\text{surj}"'] \& T\;(=|G|)\arrow[r,-,to path={([xshift=4pt,yshift=2pt]\tikztostart.east) .. controls ([xshift=-1pt,yshift=2pt]\tikztostart.east) and ([xshift=-1pt]\tikztostart.east) .. ([xshift=4pt]\tikztostart.east) -- (\tikztotarget.west) -- ([xshift=-2pt,yshift=2pt]\tikztotarget.west)}]\arrow[d,"q"',"\text{surj}"] \& L \\
I\;(=F)\arrow[r,"\varphi"] \& J\;(=G) \&
\end{tikzcd}}
où $S\subset T$ sont des sous. ens. de $L$, $\varphi:I\longrightarrow J$ une application \uline{croissante} d'ensembles ordonnés, et \deleted{$\unsure{p}$}
\[p:S-I,\quad q:T-J\]
des appl. \uline{surjectives} ($\varphi$ pas \unsure{néc.|nec.} surjective)

Alors l'inclusion $i$ est continue pour les top. \unsure{corresp|comp|can.} sur $S,T$ définies par leurs préordres respectifs (images inv. des top. sur $I,J$) L'image inverse d'une partie \inserted{$G$-}modérée de $T$ est une partie $F$-modérée de $S$, cette notion se \unsure{raméne|ramène} à la notion
\[J'\longrightarrow\varphi^{-1}(J')\]
au niveau des parties de $J,I$. L'image \deleted{\unsure{inv.}}dir. d'une partie fermée
